\documentclass[11pt]{article}
\usepackage[a4paper,margin=26mm]{geometry}
\usepackage{amsmath,amssymb,amsthm}
\usepackage[hidelinks]{hyperref}
\newtheorem{theorem}{Theorem}
\title{Parameter concavity of the generalized inverse hyperbolic sine}
\author{Research draft; independent review pending}
\date{October 28, 2026}
\begin{document}
\maketitle
\begin{abstract}
For each $x>0$, the function
$p\mapsto\int_0^x(1+t^p)^{-1/p}\,dt$ is strictly concave on
$(\tfrac12\log 2,\infty)$. This proves the parameter-concavity statement
of Conjecture~5.1 in the 2019 survey of Yin, Huang, Wang and Lin.
The proof combines an explicit second derivative with the classical
binary-entropy bound. A literal Lean certificate verifies the integral
statement using only the three standard logical axioms.
\end{abstract}

For $p>0$ and $x\geq0$, write
\begin{equation}\label{eq:def}
 A(p,x)=\operatorname{arcsinh}_p(x)
       =\int_0^x(1+t^p)^{-1/p}\,dt.
\end{equation}
Conjecture~5.1 of \cite{YHWL} asks whether $p\mapsto A(p,x)$ is strictly
concave on $(1,\infty)$ for every fixed $x\in(0,1)$. We prove a larger
parameter range and allow every $x>0$.

\begin{theorem}\label{thm:main}
For every $x>0$, the function $p\mapsto A(p,x)$ is strictly concave on
$\bigl(\tfrac12\log 2,\infty\bigr)$.
\end{theorem}

\begin{proof}
For $u>0$, define
\[
 H(u)=\log(1+u)-\frac{u\log u}{1+u}.
\]
With $r=u/(1+u)\in(0,1)$ this becomes the binary entropy
$-r\log r-(1-r)\log(1-r)$, so $H(u)>0$. Also
\[
 H'(u)=-\frac{\log u}{(1+u)^2}.
\]
Thus $H$ increases up to $u=1$ and decreases afterwards. In particular,
\begin{equation}\label{eq:entropy}
 0<H(u)\leq H(1)=\log 2.
\end{equation}

Fix $t>0$, and set $f(p,t)=(1+t^p)^{-1/p}$ and $u=t^p$.
Direct differentiation gives
\[
 \partial_p\log f(p,t)=\frac{H(u)}{p^2},\qquad
 \partial_p H(t^p)=-\frac{p(\log t)^2u}{(1+u)^2}.
\]
It follows that
\begin{equation}\label{eq:deriv}
 \partial_p^2 f(p,t)=\frac{f(p,t)}{p^4}
 \left[H(u)^2-2pH(u)-\frac{p^3(\log t)^2u}{(1+u)^2}\right].
\end{equation}
If $p>\tfrac12\log2$, then \eqref{eq:entropy} yields
$H(u)<2p$. The first two terms in the bracket equal
$H(u)(H(u)-2p)<0$, and its last term is nonpositive. Hence
$\partial_p^2 f(p,t)<0$ for every $t>0$ in the stated parameter range.

Take distinct $p_0,p_1>\tfrac12\log2$ and $0<\lambda<1$.
Strict concavity of $f(\cdot,t)$ gives, for every $t>0$,
\[
 f\bigl(\lambda p_0+(1-\lambda)p_1,t\bigr)
 >\lambda f(p_0,t)+(1-\lambda)f(p_1,t).
\]
Each integrand is continuous on $[0,x]$, including $t=0$, where its
value is one. The Jensen gap is continuous and positive on $(0,x]$;
its integral is therefore strictly positive. Integrating the last display
proves the required strict inequality for $A(\cdot,x)$.
\end{proof}

Since $\tfrac12\log2<1$, Theorem~\ref{thm:main} implies exactly the
strict-concavity statement in \cite[Conjecture~5.1]{YHWL}. Its strict
Tur\'an consequence, for $p>2$ and $0<x<1$, follows from positivity and
the arithmetic--geometric mean inequality:
\[
 A(p,x)^2>A(p-1,x)A(p+1,x).
\]
No optimality of the interval in Theorem~\ref{thm:main} is asserted.

\paragraph{Source scope.}
The 2019 question uses $(1,\infty)$ as its parameter domain.
The older \cite[Conjecture~2.4]{BBV} used $(0,\infty)$, whereas its
Theorem~2.6 proves separate parameter concavity for the two-parameter
integral $\operatorname{arcsinh}_{p,q}$ with one parameter held fixed.
The proof above differentiates the diagonal $p=q$ explicitly; separate
concavity alone does not justify that step. The binary-entropy estimate
is classical and is not presented as a new method.

\paragraph{Formal certificate and review status.}
The Lean development defines \eqref{eq:def} using real powers and a genuine
interval integral. Its closed endpoint
\texttt{ArcsinhPositive.source\_conjecture} states the exact $p>1$ source
claim. The stronger endpoint
\texttt{sourceArcsinh\_strictConcaveOn} states Theorem~\ref{thm:main} with
only $x>0$ as a hypothesis. The certificate proves the entropy identity,
both derivatives, pointwise strict concavity, continuity at $t=0$, and
strictness of the integrated Jensen inequality. It does not require
differentiation under the integral sign.

A fresh isolated replay passes all 19 declarations in five modules with
warnings treated as errors. Their axiom reports contain only
\texttt{propext}, \texttt{Classical.choice} and \texttt{Quot.sound}.
The pinned toolchain is Lean~4.34.1, with mathlib revision\\
\texttt{d13f23b723b8a846827a245b89c10fc7d3f11612}.

This is a research draft. Independent source, mathematical, formal,
novelty and substantive review remain pending; no discovery acceptance
or worldwide priority is claimed. In particular, possible subsumption by
a generic reciprocal-norm parameter theorem remains an explicit review
question. The entire result is one proposed family; supporting identities
and consequences are not separate outcomes.

\begin{thebibliography}{9}
\bibitem{YHWL} L. Yin, L.-G. Huang, Y.-L. Wang and X.-L. Lin,
\emph{A survey for generalized trigonometric and hyperbolic functions},
J. Math. Inequal. \textbf{13}(3) (2019), 833--854.
\href{https://doi.org/10.7153/jmi-2019-13-58}{doi:10.7153/jmi-2019-13-58}.
\bibitem{BBV} \'A. Baricz, B. A. Bhayo and M. Vuorinen,
\emph{Tur\'an type inequalities for generalized inverse trigonometric functions},
Filomat \textbf{29}(2) (2015), 303--313.
\href{https://doi.org/10.2298/FIL1502303B}{doi:10.2298/FIL1502303B}.
\end{thebibliography}
\end{document}
